\begin{helper}[Exact edges selected from a defective crossing]
Let $d$ be a defective color in a reduced Scheide configuration, put
$i=\gammaIndex(d)$, $X=V(T)$,
$U=\{u^1_{\gamma_i},u^2_{\gamma_i}\}$, and $Z=X\setminus U$.
Suppose that an exact edge identity $g_1$ has color $d$ and is incident
with a vertex $x\in U$ and a vertex $z_1\in Z$.  Then the two vertices of
$U$ can be named $u_1,u_2$, with $x=u_1$, and there are exact edge
identities $g_2,g_3$ and vertices $z_3,w_3$ such that
\[
 \begin{gathered}
 u_1\ne u_2,\quad u_1,u_2\in X,\quad
 v^0_d,y_0\in Z,\\
 \alpha_0\ne\gamma_i,\quad \alpha_0\ne d,\quad \gamma_i\ne d,\\
 \operatorname{ends}(g_1)=\{u_1,z_1\},\qquad
 \varphi(g_2)=\alpha_0,\quad
 \operatorname{ends}(g_2)=\{u_1,u_2\},\\
 g_3\in\partial_d(X),\quad g_3\ne f^0_d,\quad
 z_3\in Z,\quad w_3\notin X,\quad
 \operatorname{ends}(g_3)=\{z_3,w_3\}.
 \end{gathered}
\]
Moreover $g_1,g_2,g_3$ are pairwise distinct.
\end{helper}
\begin{proof}
Write the two members of $U$ as $u_1=x$ and $u_2$ for the other member.
The two block positions in the ordered \noderef{TreeSequence} are
distinct, so $u_1\ne u_2$; the structural position equations in
\noderef{ScheideReducedConfiguration} put both
vertices in $X$.  The two stated incidences of $g_1$, together with the
endpoint membership in \noderef{Incident} and the loopless law in
\noderef{LooplessMultigraph}, give
$\operatorname{ends}(g_1)=\{u_1,z_1\}$.

We locate the two distinguished base vertices.  By
\noderef{ScheideGlobalMaximumOrder}, $X$ is elementary.  Membership of
$v^0_d$ in the exit set, supplied by \noderef{ScheideDefectiveSetup},
and the exit-covering clause give a block color missing at $v^0_d$.
The same block color is missing at a base vertex by the block-missing
clause.  Elementariness forces those two vertices to coincide.  Hence
$v^0_d$ is a base vertex.  The first vertex $y_0$ is also a base vertex.
Vertex injectivity and the structural position equations separate every
base position from both positions of the $i$-block.  Thus
$v^0_d,y_0\in Z$.

The setup \noderef{ScheideDefectiveSetup} already states
$\alpha_0\ne\gamma_i$.  Since $d$ is defective,
\noderef{ScheideGlobalMaximumBoundaryFacts} gives at least three exact
$d$-colored boundary identities.  Choose one.  Closedness from
\noderef{ScheideGlobalMaximumOrder} says that the color of a boundary
edge is missing at no vertex of $X$.  By \noderef{MissingColors},
therefore, $d$ is present at every
vertex of $X$.  It is distinct from $\alpha_0$, which is missing at
$y_0$, and from $\gamma_i$, which is missing at $v^0_d$.

Now $\alpha_0$ is missing at the vertex $y_0\in Z$ and is distinct from
$\gamma_i$.  Applying \noderef{ScheideStructuralInternalPair} gives an
exact identity $g_2$ of color $\alpha_0$ with endpoint pair
$\{u_1,u_2\}$.

The setup gives
$\varphi(f^0_d)=d$, endpoints $\{v^0_d,u^0_d\}$, and
$u^0_d\notin X$.  Hence $f^0_d\in\partial_d(X)$ and its endpoint in
$X$ lies in $Z$.  Properness in the globally maximum field of
\noderef{ScheideReducedConfiguration}, interpreted by
\noderef{IsProperEdgeColoring}, permits only one $d$-identity incident
with either $u_1$ or $u_2$.  The internal identity $g_1$ already uses the
$d$-incidence at $u_1$.  Consequently at most one member of
$\partial_d(X)$ has its endpoint in $U$.  Since this boundary-color set
has at least three members and $f^0_d$ has its endpoint in $Z$, there is
a member $g_3\ne f^0_d$ whose endpoint in $X$ lies in $Z$.  Unfolding
\noderef{BoundaryColorEdges} and \noderef{BoundaryEdges} gives vertices
$z_3\in Z$ and $w_3\notin X$ with the displayed endpoint equation.

Finally, $g_1$ is internal to $X$ whereas $g_3$ is a boundary identity,
so they differ.  The identity $g_2$ has color $\alpha_0$, while $g_1$
and $g_3$ have color $d\ne\alpha_0$; hence it differs from both.  These
arguments count exact edge identities and therefore remain valid in the
presence of parallel edges.
\end{proof}
